\documentclass{beamer}
%\usepackage[utf8]{inputenc}
\usepackage[T1]{fontenc}
\usepackage{mathabx}
\usepackage{eulervm}
\usepackage{natbib}

\usepackage{fontspec}
%\setmainfont{EBGaramond12-Regular}
\setmainfont{MinionPro-Regular}
%\setmainfont{LibertinusSerif-Regular}

\usepackage[scheme=plain]{ctex}
\setCJKmainfont[BoldFont = SourceHanSerifSC-Bold,
        ItalicFont = KaiTi]
        {SimSun}

%% Load the markdown package
\usepackage[citations,footnotes,definitionLists,hashEnumerators,smartEllipses,tightLists=false,pipeTables,tableCaptions,hybrid]{markdown}
%%begin novalidate
\markdownSetup{rendererPrototypes={
 link = {\href{#2}{#1}},
 headingOne = {\section{#1}},
 headingTwo = {\subsection{#1}},
 headingThree = {\begin{frame}\frametitle{#1}},
 headingFour = {\begin{block}{#1}},
 horizontalRule = {\end{block}}
}}
%%end novalidate

\usetheme{metropolis}
\usefonttheme{serif}
\usecolortheme{rose}


\title{Writing in Markdown, Presenting with \LaTeX}
\author{TomBen}
\institute{中文支持}
\titlegraphic{\hfill\includegraphics[height=1.5cm]{logo.pdf}}

\begin{document}

\maketitle

\frame{\tableofcontents}

\begin{markdown}
%%begin novalidate

# How?

### We can do that?

- Yeah to some extent, with \texttt{markdown} package :-)
    - __$\hash$__ and __$\hash\hash$__ for section and subsection headers (in ToC)
    - Redefine __$\hash\hash\hash$__ to start a frame and frametitle
    - (Nested) bullet and numbered lists
    - Text formatting (*italic*, **bold becomes italic + alerted**) 
    - Redefine __$\hash\hash\hash\hash$__ to start a block with title \linebreak
      and __\texttt{-{}-{}-{}-}__ to end the block
    - ___Compile with \texttt{-{}-shell-escape}___ (Overleaf does this already)
- (Alternative approaches: Pandoc, wikitobeamer)

\end{frame}

%%%%%%%%%%%%%%%%%%%%%%

### Caveats

- Nothing too complicated! 
- No verbatim or fragile stuff!
- No $\hash$ and \textunderscore{} characters!\linebreak 
  (I used `$\hash$` and `\textunderscore`)
- Can't pass options to frames
- __Need to write \texttt{\textbackslash end\string{frame\string}} manually!__

\end{frame}

%%%%%%%%%%%%%%%%%%%%%%


%%% # and ## can still be used as sections and subsections if you prefer
# Example

## Proposed Menus

%%% ### starts a frame + frametitle
### Breakfast Menu

%%% bulleted lists as usual 

- Eggs
    * scrambled
    * sunny-side-up
- Coffee
    * Americano
    * Long black
- Tea
    * Darjeeling
    * English Breakfast

%%% Due to the complicatedness of beamer frames, \end{frame} MUST appear in the source code itself and cannot be "hidden" in another command

\end{frame}

%%%%%%%%%%%

### Lunch Menu

- Spaghetti
    * Bolognese
    * Aglio olio
- Sandwiches
    * Egg
    * Ham
    * Tuna

\end{frame}

%%%%%%%%%%%

## Budgeting

### Projected Profit

1. And the answer is...
2. $f(x)=\sum_{n=0}^\infty\frac{f^{(n)}(a)}{n!}(x-a)^n$
    #. How do we _know_ that?
    #. __Maths!__

$$\int f(x)\mathrm{d}x=\lim_{0 \to + \infty}g(x)+\lambda$$
\end{frame}

### Testing blocks

#### This is a block!

- Here is some content.
- Here's more contents.

---

\end{frame}


### Citations

- This is a citation [@Knuth92]
- Works like `natbib` syntax [see @Er01 p.26]
- A citation [@Simpson]

\end{frame}


### Pipe Tables

- Use `pipeTables` and `tableCaptions` options
- Available since `markdown` v2.8.0

| Right | Left | Default | Center |
|------:|:-----|---------|:------:| 
|  12   |  12  |  12     |   12   | 
| 123   |  123 |   123   |  123   | 
|   1   |    1 |     1   |    1   | 

  : Demonstration of pipe table syntax.
  
\end{frame}

%%novalidate
\end{markdown}

\begin{frame}
\renewcommand{\bibfont}{\footnotesize}
\frametitle{Bibliography}

\bibliographystyle{apalike}
\bibliography{refs}

\end{frame}


\end{document}
